\section*{Hoofdstuk \ref{ch:g2:measuring-length} --- Lengte meten}

\begin{solution}{exo:g2:measuring-length:1}
Dichtbij: rug aan rug gaan staan. Ver weg: jezelf in
\omterm{def:g2:measuring-length:units}{centimeters} \omterm{def:g2:measuring-length:measure}{meten} en het getal opsturen --- een getal kan in een
brief reizen, een vergelijking rug aan rug niet.
\end{solution}

\begin{solution}{exo:g2:measuring-length:2}
Geen van beiden heeft verkeerd geteld. De stappen van Mara zijn
korter, dus passen er meer in hetzelfde lokaal. Verschillende
eenheden geven verschillende getallen voor dezelfde lengte.
\end{solution}

\begin{solution}{exo:g2:measuring-length:3}
De \omterm{def:g2:measuring-length:units}{centimeter} is voor iedereen hetzelfde --- hij is van niemands
lichaam. Een meting in \omterm{def:g2:measuring-length:units}{centimeters} kan door ieder ander met elke
liniaal worden nagekeken.
\end{solution}

\begin{solution}{exo:g2:measuring-length:4}
De mier: \unit{cm} (zelfs minder!); de deur: \unit{m}; het
schoolplein: \unit{m}; de breedte van dit boek: \unit{cm}.
\end{solution}

\begin{solution}{exo:g2:measuring-length:5}
Nee --- tussen de rand van een liniaal en het nulstreepje zit meestal
een klein stukje leeg, dus het potlood is iets korter dan $13$
\omterm{def:g2:measuring-length:units}{centimeter}. Tim vergat het \emph{nulstreepje}, en niet de rand, bij
het uiteinde van het potlood te leggen.
\end{solution}

\begin{solution}{exo:g2:measuring-length:6}
De antwoorden verschillen; gebruikelijk: hand ongeveer \qty{7}{cm}
breed, lepel ongeveer \qty{15}{cm}, beker ongeveer \qty{9}{cm}. Wat
telt: elk getal draagt zijn \unit{cm}.
\end{solution}

\begin{solution}{exo:g2:measuring-length:7}
$100 + 100 = 200$: de deur is \qty{200}{cm} hoog.
\end{solution}

\begin{solution}{exo:g2:measuring-length:8}
De antwoorden verschillen; de meeste bedden komen in de buurt van
\qty{2}{m}. Het gaat erom de gok met het gemeten getal te
vergelijken.
\end{solution}

\begin{solution}{exo:g2:measuring-length:9}
Het lint van \qty{45}{cm} is langer: $45 - 38 = 7$, dus \qty{7}{cm}
meer. Zonder getallen vind je de langste ook door ze naast elkaar te
leggen --- maar ``\qty{45}{cm}'' past in een brief, en ``naast
elkaar'' niet.
\end{solution}

\begin{solution}{exo:g2:measuring-length:10}
Niemand anders heeft opa's benen: een klein kind dat dezelfde tuin
afstapt telt misschien $60$ kleine passen, en een bezoeker kan
``$30$ passen'' niet nakijken zonder opa erbij. Gemeten in
\omterm{def:g2:measuring-length:units}{meters} --- voor iedereen hetzelfde --- kun je de lengte van de tuin
opschrijven, versturen en door iedereen laten nakijken.
\end{solution}
